\documentclass[10pt]{article}
\usepackage[T1]{fontenc}
\usepackage[utf8]{inputenc}
\usepackage[french]{babel}
\usepackage{amsmath,amssymb,mathrsfs,mathpazo}
\usepackage{geometry}
\geometry{a4paper,margin=2cm}
\usepackage{graphicx,array,tabularx,tabularray,colortbl,multirow,multicol,booktabs,enumitem,pifont,ulem,fancybox,diagbox,cancel}
\usepackage{pstricks,pstricks-add,pst-plot,pst-tree,pst-node,pst-text,pst-func,pst-math,pst-bspline,pst-3dplot,pst-eucl,pst-solides3d}
\let\pscalculate\undefined
\usepackage{pst-bezier}
\usepackage{tikz,pgfplots,tkz-tab}
\usetikzlibrary{arrows,arrows.meta,patterns,calc,positioning,trees,decorations.markings,shapes.geometric,shapes.misc,shapes.arrows,intersections,decorations.pathmorphing,decorations.pathreplacing,backgrounds,shadows,babel}
\pgfplotsset{compat=1.18}
\usepgfplotslibrary{fillbetween}
\usepackage[np]{numprint}
\usepackage{listings,makecell,pas-tableur,twoopt}
\let\barre\undefined
\usepackage[pstricks]{tablvar}
\definecolor{darkgreen}{RGB}{0,100,0}
\providecommand{\red}{\color{red}}
\providecommand{\blue}{\color{blue}}
\providecommand{\green}{\color{green}}
\providecommand{\black}{\color{black}}
\providecommand{\white}{\color{white}}
\providecommand{\gray}{\color{gray}}
\providecommand{\cyan}{\color{cyan}}
\providecommand{\magenta}{\color{magenta}}
\providecommand{\yellow}{\color{yellow}}
\providecommand{\violet}{\color{violet}}
\providecommand{\orange}{\color{orange}}
\providecommand{\euro}{\texteuro}
\providecommand{\EUR}{\texteuro}
\providecommand{\ieme}{\textsuperscript{e}}
\providecommand{\degres}{\ensuremath{^\circ}}
\providecommand{\degree}{\ensuremath{^\circ}}
\providecommand{\vv}[1]{\overrightarrow{#1}}
\providecommand{\Coord}[1]{#1}
\providecommand{\decosix}{}
\providecommand{\starredbullet}{\ensuremath{\bullet}}
\pagestyle{empty}
\setlength{\parindent}{0pt}
\usepackage{adjustbox}\DeclareUnicodeCharacter{2236}{:}\DeclareUnicodeCharacter{2190}{\ensuremath{\leftarrow}}
\providecommand{\R}{\mathbb{R}}
\providecommand{\N}{\mathbb{N}}
\providecommand{\Z}{\mathbb{Z}}
\providecommand{\Q}{\mathbb{Q}}
\providecommand{\D}{\mathbb{D}}
\providecommand{\C}{\mathbb{C}}
\providecommand{\e}{{\mathrm{e}}}
\providecommand{\vect}[1]{\overrightarrow{#1}}
\providecommand{\vectt}[1]{\overrightarrow{\mathrm{#1}}}
\providecommand{\vvt}[1]{\overrightarrow{\mathrm{#1}}}
\providecommand{\np}[2][]{{#2~#1}}
\providecommand{\vv}[1]{{\overrightarrow{#1}}}
\renewcommand{\d}{{\mathrm{d}}}
\providecommand{\degres}{{^\circ}}
\providecommand{\degre}{{^\circ}}
\providecommand{\blue}{\color{blue}}
\providecommand{\red}{\color{red}}
\providecommand{\enskip}{\quad}
\providecommand{\llbracket}{\lbrack\!\lbrack}
\providecommand{\rrbracket}{\rbrack\!\rbrack}
\providecommand{\euro}{{\text{€}}}
\providecommand{\no}{{\text{n°}}}
\providecommand{\interval}[2]{{[#1;#2]}}
\providecommand{\intervFO}[2]{{[#1;#2[}}
\providecommand{\dotfill}{\ldots}

\providecommand{\ds}{\displaystyle}
\providecommand{\textding}[1]{#1}
\providecommand{\Oij}{\(\mathrm{O};\vec{\imath},\vec{\jmath}\)}
\providecommand{\Oijk}{\(\mathrm{O};\vec{\imath},\vec{\jmath},\vec{k}\)}
\providecommand{\Ouv}{\(\mathrm{O};\vec{u},\vec{v}\)}
\providecommand{\intervFF}[2]{\ensuremath{[#1;#2]}}
\renewcommand{\euro}{\ifmmode\text{\texteuro}\else\texteuro\fi}
\DeclareUnicodeCharacter{00B0}{\ifmmode^{\circ}\else\textdegree\fi}
\DeclareUnicodeCharacter{20AC}{\ifmmode\text{\texteuro}\else\texteuro\fi}
% Énoncés : Éducation nationale. Transcription et corrigés : APMEP, Denis Vergès et contributeurs. Découpage et classement : CoopMaths, Sébastien Lozano, Éric Elter et contributeurs. Adaptation de la présentation : Maths Explorer.
% CC BY-SA : https://coopmaths.fr/www/about/
% https://coopmaths.fr/alea/static/eam/2026/tex/eam_technologique_2026_09_metropole_1_cor.tex
\begin{document}
\noindent\textit{APMEP / CoopMaths. CC BY-SA. Présentation adaptée par Maths Explorer.}\par\medskip
\emph{Groupe témoin de \np{1000} bovins : 200 malades (180 testés positifs, 20 négatifs) et
800 non malades (45 positifs, 755 négatifs) ; au total 225 tests positifs.
On note $M$ l'évènement \og le bovin est malade \fg{} et $T$ l'évènement \og le test est positif \fg.}

\begin{enumerate}
	\item Le bovin est choisi au hasard : on est en situation d'équiprobabilité.
	\[P(M) = \dfrac{200}{\np{1000}} = 0,2.\]

	\item $M \cap T$ : \og le bovin est malade et son test est positif \fg.
	\[P(M \cap T) = \dfrac{180}{\np{1000}} = 0,18.\]
	La probabilité que le bovin soit malade et que son test soit positif est égale à $0,18$.

	\item Parmi les 200 bovins malades, 180 ont un test positif :
	\[P_M(T) = \dfrac{180}{200} = 0,9.\]

	\item \og Ne pas être atteint sachant que le test est positif \fg{} se traduit par
	$P_T\left(\overline{M}\right)$. Parmi les 225 bovins testés positifs, 45 ne sont pas malades :
	\[P_T\left(\overline{M}\right) = \dfrac{45}{225} = \dfrac{45}{5 \times 45} = \dfrac{1}{5} = 0,2.\]
	Or $0,2 > 0,15$ : le test \textbf{n'est pas fiable}.

	\item
	\begin{enumerate}
		\item Le tirage est assimilé à un tirage avec remise : à chaque niveau,
		$P(M) = 0,2$ et $P\left(\overline{M}\right) = 1 - 0,2 = 0,8$.

		\begin{center}
			\begin{tikzpicture}[x=2.8cm, y=0.9cm]
				
				\node (root) at (0,0) {};
				\node (M) at (1,1.2) {$M$};
				\node (Mb) at (1,-1.2) {$\overline{M}$};
				\node (M1) at (2,1.9) {$M$};
				\node (M1b) at (2,0.5) {$\overline{M}$};
				\node (M2) at (2,-0.5) {$M$};
				\node (M2b) at (2,-1.9) {$\overline{M}$};
				
				\draw (root.east) -- (M) node[pos=0.5,above] {0,2};
				\draw (root.east) -- (Mb) node[pos=0.5,below] {0,8};
				
				\draw (M.east) -- (M1) node[pos=0.5,above] {0,2};
				\draw (M.east) -- (M1b) node[pos=0.5,below] {0,8};
				\draw (Mb.east) -- (M2) node[pos=0.5,above] {0,2};
				\draw (Mb.east) -- (M2b) node[pos=0.5,below] {0,8};
				
				\node at (1,2.6) {\small 1\up{er} bovin};
				\node at (2,2.6) {\small 2\up{e} bovin};
			\end{tikzpicture}
		\end{center}

		\item L'évènement $A$ correspond au chemin $M$ -- $M$ de l'arbre :
		\[P(A) = 0,2 \times 0,2 = 0,04.\]
		La probabilité que les deux bovins choisis soient malades est égale à $0,04$.
	\end{enumerate}
\end{enumerate}

\end{document}
